\subsection{Bài tập bài 2: Giải các bài toán QHTT bằng phương pháp đơn hình}

% --- Câu a ---
\begin{exerciseblock}{Bài tập a}
\begin{alignat*}{4}
    f(x) = \quad & 3x_1 & {}+{} & 2x_2 & {}+{} & 5x_3 & \rightarrow \max \\
    \text{s.t.} \quad & x_1 & {}+{} & x_2 & {}+{} & 2x_3 & \le 4000 \\
    & x_1 & {}+{} & x_2 & {}+{} & 3x_3 & \le 9000 \\
    & 3x_1 & {}+{} & x_2 & {}+{} & 2x_3 & \le 8000 \\
    & \rlap{$x_j \ge 0, \quad j = 1,2,3$}
\end{alignat*}
\end{exerciseblock}

% --- Câu b ---
\begin{exerciseblock}{Bài tập b}
\begin{alignat*}{4}
    f(x) = \quad & 4x_1 & {}+{} & 2x_2 & {}+{} & 6x_3 & \rightarrow \max \\
    \text{s.t.} \quad & x_1 & {}+{} & x_2 & {}+{} & 2x_3 & \le 8000 \\
    & x_1 & {}+{} & x_2 & {}+{} & 3x_3 & \le 16000 \\
    & 3x_1 & {}+{} & x_2 & {}+{} & 2x_3 & \le 12000 \\
    & \rlap{$x_j \ge 0, \quad j = 1,2,3$}
\end{alignat*}
\end{exerciseblock}

% --- Câu c ---
\begin{exerciseblock}{Bài tập c}
\begin{alignat*}{4}
    f(x) = \quad & -2x_1 & {}+{} & 5x_2 & {}+{} & 3x_3 & \rightarrow \max \\
    \text{s.t.} \quad & x_1 & {}+{} & 2x_2 & {}+{} & 3x_3 & \le 8000 \\
    & x_1 & {}+{} & 3x_2 & {}+{} & x_3 & \le 10000 \\
    & x_1 & {}+{} & 2x_2 & {}+{} & x_3 & \le 4000 \\
    & \rlap{$x_j \ge 0, \quad j = 1,2,3$}
\end{alignat*}
\end{exerciseblock}

% --- Câu e ---
\begin{exerciseblock}{Bài tập e}
\begin{alignat*}{4}
    f(x) = \quad & 4x_1 & {}-{} & 4x_2 & {}+{} & 6x_3 & \rightarrow \max \\
    \text{s.t.} \quad & x_1 & {}+{} & x_2 & {}+{} & 3x_3 & \le 21000 \\
    & 3x_1 & {}+{} & x_2 & {}+{} & 2x_3 & \le 12000 \\
    & x_1 & {}+{} & x_2 & {}+{} & 2x_3 & \le 8000 \\
    & \rlap{$x_j \ge 0, \quad j = 1,2,3$}
\end{alignat*}
\end{exerciseblock}

% --- Câu g (Biến có điều kiện xi >= 1) ---
\begin{exerciseblock}{Bài tập g}
\begin{alignat*}{4}
    f(x) = \quad & 6x_1 & {}+{} & 3x_2 & {}-{} & 7x_3 & \rightarrow \min \\
    \text{s.t.} \quad & -x_1 & {}-{} & x_2 & {}+{} & 3x_3 & \le 3 \\
    & -2x_1 & {}-{} & 4x_2 & {}+{} & 2x_3 & \le 9 \\
    & 4x_1 & {}+{} & 2x_2 & {}-{} & 3x_3 & \le 2 \\
    & \rlap{$x_j \ge 1, \quad j = 1,2,3$}
\end{alignat*}
\end{exerciseblock}

% --- Câu h ---
\begin{exerciseblock}{Bài tập h}
\begin{alignat*}{4}
    f(x) = \quad & -5x_1 & {}+{} & 3x_2 & {}+{} & 3x_3 & \rightarrow \min \\
    \text{s.t.} \quad & 3x_1 & {}+{} & x_2 & {}+{} & x_3 & \le 11 \\
    & 2x_1 & {}+{} & 3x_2 & {}+{} & x_3 & \le 8 \\
    & 2x_1 & {}+{} & x_2 & {}+{} & 4x_3 & \le 4 \\
    & \rlap{$x_j \ge 0, \quad j = 1,2,3$}
\end{alignat*}
\end{exerciseblock}

% --- Câu k (Hỗn hợp đẳng thức và bất đẳng thức) ---
\begin{exerciseblock}{Bài tập k}
\begin{alignat*}{5}
    f(x) = \quad & -2x_1 & {}+{} & x_2 & {}+{} & x_4 & & \rightarrow \min \\
    \text{s.t.} \quad & x_1 & {}+{} & x_2 & {}-{} & x_3 & & \le 15 \\
    & x_1 & {}+{} & x_2 & {}+{} & x_3 & {}+{} x_4 & = 27 \\
    & 2x_1 & {}-{} & x_2 & {}-{} & x_3 & & \le 18 \\
    & \rlap{$x_j \ge 0, \quad j = \overline{1,4}$}
\end{alignat*}
\end{exerciseblock}

% --- Câu p (Dạng chính tắc có sẵn biến cơ sở) ---
\begin{exerciseblock}{Bài tập p}
\begin{alignat*}{4}
    f(x) = \quad & 6x_1 & {}+{} & 10x_2 & {}+{} & 12x_3 & \rightarrow \max \\
    \text{s.t.} \quad & x_1 & {}+{} & 2x_2 & {}+{} & 4x_3 & \le 32 \\
    & 2x_1 & {}+{} & x_2 & {}+{} & 2x_3 & \le 42 \\
    & 3x_1 & {}+{} & 2x_2 & {}+{} & 2x_3 & = 30 \\
    & \rlap{$x_j \ge 0, \quad j = 1,2,3$}
\end{alignat*}
\end{exerciseblock}